\documentclass{article}
\usepackage[margin=1in]{geometry}
\usepackage{amsmath}
\usepackage{amssymb}

\begin{document}

\begin{center}
{\Large\bfseries Bearings Challenge --- Answer Key}
\end{center}

\section*{Introduction}
\begin{enumerate}
\item Gold: \(100\) km from the start on bearing \(123.5^\circ\), i.e.\ \(83.38868\) km east and \(55.19370\) km south.
\item Bob: due south \(55.19370\) km, then due east \(83.38868\) km.
\item Bob total distance: \(138.58\) km; time \(=6.93\) days.
\item Rob: \(78.0557\) km on bearing \(135^\circ\), then east. Total \(=106.35\) km; time \(=5.32\) days; Rob arrives about \(1\) day \(15\) hours before Bob.
\item Best two-leg route: \(50.00190\) km on bearing \(124^\circ\), then \(50.00190\) km on bearing \(123^\circ\); this takes less than a minute more than \(5\) days, so even with a \(7\)-hour delay it beats Rob's time.
\end{enumerate}

\section*{Challenge I: Integer Distances}
\begin{enumerate}
\item
\begin{enumerate}
\item Bob, after rounding legs to whole km: \(434\) m from gold.
\item Rob, after rounding legs to whole km: \(237\) m from gold.
\end{enumerate}

\item Rounded quickest journey: \(50\) km on \(124^\circ\), then \(50\) km on \(123^\circ\); endpoint is \(99.99619\) km from start and \(3.81\) m from gold, so not within \(1\) m; gold remains buried at this point.

\item Four-leg displacement: \(0.5316\) cm east and \(91.3783\) cm north.

\item Same route rotated \(90^\circ\) anticlockwise: \(0.5316\) cm north and \(91.3783\) cm west, i.e.\ \(-91.3783\) cm east.

\item Use \(A\approx -59868.41\), \(B\approx -91604.77\); equivalently, \(59868\) reversed copies of the first 4-leg route and \(91605\) reversed copies of the rotated route. Total distance \(=908838\) km. This ends about \(11\) cm east and \(37\) cm north of the gold, about \(39\) cm away, so the gold is found.

\item Yes. For any target, real \(A,B\) can solve the two equations; rounding the repetition counts to integers leaves less than \(50\) cm error in each coordinate, hence the endpoint is within \(1\) m.
\end{enumerate}

\section*{Challenge II: Two Legs Only}
\begin{enumerate}
\item For the considered long first leg, the only return bearing that can get within \(100\) km of the start is the exact reverse bearing, e.g.\ \(270^\circ\) after \(090^\circ\); by symmetry, the only candidate return bearing is
\[
XYZ^\circ=(ABC^\circ+180^\circ)\bmod 360^\circ.
\]

\item Necessary condition:
\[
|M-N|\le 100,
\]
so \(N\) must lie in \([M-100,M+100]\).

\item Upper bound:
\[
360\cdot6000\cdot360\cdot203 \approx 1.58\times10^{11}
\]
possible destinations within \(101\) km.

\item Not necessarily, but possibly; checking all finite possibilities would be required.

\item Upper bound counts:
\begin{enumerate}
\item \(M\le100\): \(2.66\times10^9\).
\item \(100<M\le500\): \(5.26\times10^9\) (at most \(180\) return bearings).
\item \(500<M\le2000\): \(3.18\times10^9\) (at most \(29\) return bearings).
\item \(2000<M\le6000\): \(2.05\times10^9\) (at most \(7\) return bearings).
\end{enumerate}
Total \(\le 1.32\times10^{10}\).

\item Reversing the leg order gives the same endpoint, so the bound halves:
\[
U\le 6.6\times10^9<10^{10}.
\]

\item Not every location is reachable within \(1\) m: the reachable area is at most \(6.6\times10^9\pi\ \text{m}^2\), less than the circle area \(10^{10}\pi\ \text{m}^2\); at least \(3\times10^9\pi\ \text{m}^2\) is unreachable.

\item Naive brute force would take over \(42\) hours; the vectorised NumPy version takes about \(40\) minutes and finds \(82\) candidate two-leg journeys ending within \(15\) m of the gold. Plotting them shows none is within \(1\) m, so the gold cannot be found by a two-leg journey.

\item One two-leg journey ends within \(2\) m of the gold: \(4937\) km on bearing \(004^\circ\), then \(4987\) km on bearing \(183^\circ\); final coordinates \((83.38830,\,-55.19177)\).
\end{enumerate}

\end{document}